\documentclass[a4paper,10pt]{article}
\usepackage[utf8]{inputenc}
\usepackage[T1]{fontenc}
\usepackage[empty]{fullpage}
\usepackage{xcolor}
\usepackage{hyperref}
\usepackage{geometry}
\usepackage{enumitem}
\usepackage{titlesec}
\usepackage{tabularx}

% ----------  Page margins  ----------
\geometry{left=0.5in, top=0.4in, right=0.5in, bottom=0.4in}

% ----------  Colors  ----------
\definecolor{accent}{HTML}{1A2A52}   % deep navy
\definecolor{muted}{HTML}{555555}

% ----------  Hyperlinks  ----------
\hypersetup{
	colorlinks=true,
	urlcolor=accent,
	linkcolor=accent
}

% ----------  Section formatting (Option C: uppercase, letterspaced, thin rule) ----------
\titleformat{\section}{
	\vspace{-4pt}\scshape\raggedright\large\bfseries\color{accent}
}{}{0em}{}[\color{accent}\titlerule \vspace{-4pt}]
\titlespacing*{\section}{0pt}{12pt}{6pt}

% ----------  Custom commands  ----------
\newcommand{\resumeItem}[1]{
	\item\small{#1}
	\vspace{1pt}
}

\newcommand{\resumeSubheading}[4]{
	\vspace{4pt}\item
	\begin{tabular*}{0.99\textwidth}[t]{l@{\extracolsep{\fill}}r}
		\textbf{#1} & #2 \\
		\textit{\small#3} & \textit{\small #4} \\
	\end{tabular*}\vspace{3pt}
}

\newcommand{\resumeProjectHeading}[2]{
	\vspace{5pt}\item
	\begin{tabular*}{0.99\textwidth}{l@{\extracolsep{\fill}}r}
		\small#1 & \small#2 \\
	\end{tabular*}\vspace{3pt}
}

% Indented project description (sits cleanly under the project title)
\newcommand{\resumeItemDesc}[1]{
	\begin{itemize}[leftmargin=0.32in, label={}, topsep=0pt, partopsep=0pt]
		\item \small{#1}
	\end{itemize}
	\vspace{1pt}
}

\setlist[itemize]{topsep=4pt, partopsep=0pt, itemsep=4pt, parsep=0pt}
\renewcommand{\labelitemii}{$\circ$}
\linespread{1.05}

\begin{document}
	
	%==================  HEADER (Option C: name left, contacts right)  ==================
	\begin{tabular*}{\textwidth}{@{\extracolsep{\fill}} l r}
		{\Huge \textbf{Rakshith Ajetaravi}} &
		\small \href{mailto:arakshith@iisc.ac.in}{arakshith@iisc.ac.in} \\
		{\color{accent}\small Robotics $\cdot$ Control $\cdot$ Autonomous Systems} &
		\small +91 63637 59090 \\
		& \small \href{https://github.com/rakshith-6}{GitHub} \quad $|$ \quad \href{https://www.linkedin.com/in/rakshith-ajetaravi}{LinkedIn} \\
	\end{tabular*}
	\vspace{2pt}
	\hrule height 0.6pt
	\vspace{-2pt}
	
	%==================  EDUCATION  ==================
	\section{Education}
	\begin{itemize}[leftmargin=0.15in, label={}]
		\resumeSubheading
		{Indian Institute of Science (IISc)}{Bengaluru, India}
		{M.Tech in Robotics and Autonomous Systems (CPS Department)}{Present}
		\begin{itemize}
			\item \textbf{Current Performance: 7.7 CGPA}
			\item \textbf{Relevant Courses:}
			\begin{itemize}
				\item \textbf{Foundations:}
				Foundations of Robotics; Mathematical Techniques for Robotics Systems
				\item \textbf{Control Systems (Specialization):}
				Applied Linear and Non-Linear Control; Non-linear Systems and Control; Sliding Mode Control and its Applications; Applied Optimal Control and State Estimation
				\item \textbf{Robotics Systems:}
				Swarm Robotic Systems; Motion Planning for Autonomous Systems
				\item \textbf{Machine Learning \& AI:}
				Machine Learning for Cyber-Physical Systems; Reinforcement Learning; Generative and Agentic AI in Practice
				\item \textbf{Embedded Systems:}
				Design of Cyber-Physical Systems
			\end{itemize}
			\item \textbf{Teaching Assistant:}
			\begin{itemize}
				\item E1241: Dynamics of Linear Systems (August term 2025)
				\item CP230: Motion Planning for Autonomous Systems (January term 2026)
			\end{itemize}
		\end{itemize}
		
		\resumeSubheading
		{National Institute of Engineering (NIE)}{Mysuru, India}
		{Bachelor of Engineering in Mechanical Engineering}{2022}
		\begin{itemize}
			\item \textbf{CGPA: 8.75 / 10}
			\item \textbf{Focus Areas:} Applied Mechanics and Design; Thermal and Fluids; Controls, Robotics and Mechatronics; Manufacturing, Materials and Industrial Engineering.
		\end{itemize}
		
		\resumeSubheading
		{Gopalaswamy Shishuvihara PU College (Sankalpa)}{Mysuru, India}
		{Class 12th (Karnataka State Board) -- PCMB}{2018}
		\begin{itemize}
			\item \textbf{Percentage: 81.67\%}
		\end{itemize}
		
		\resumeSubheading
		{JSS Public School, JP Nagar}{Mysuru, India}
		{Class 10th (CBSE Board)}{2016}
		\begin{itemize}
			\item \textbf{CGPA: 9.8 / 10}
		\end{itemize}
	\end{itemize}
	
	%==================  TECHNICAL SKILLS  ==================
	\section{Technical Skills}
	\begin{itemize}[leftmargin=0.15in, label={}]
		\vspace{6pt}
		\small{\item{
				\textbf{Robotics \& Simulation}{: ROS2, Gazebo, Rviz, MuJoCo, NVIDIA ISAAC Sim} \\
				\textbf{Programming}{: Python, C	, MATLAB, Julia} \\
				\textbf{Data Science \& ML}{: TensorFlow, Keras, PyTorch, Scikit-learn, Pandas, NumPy, Gymnasium, OpenCV} \\
				\textbf{Control \& Estimation}{: LQR, MPC, Sliding Mode Control, ADRC, Kalman Filtering, Reinforcement Learning} \\
				\textbf{Engineering Design}{: SolidWorks, Fusion 360, CATIA, Ansys, Qblade, Ultimaker Cura}
		}}
	\end{itemize}
	
	%==================  PROJECTS (exact PDF order, verbatim)  ==================
	\section{Academic \& Personal Projects}
	\begin{itemize}[leftmargin=0.15in, label={}]
		
		\resumeProjectHeading
		{\textbf{E-Gokart and Formula Student Vehicle} $|$ \emph{Team Captain}}{Oct'20 -- Aug'22}
		\vspace{-6pt}
		\resumeItemDesc{Involved in \textbf{design and manufacturing of electric go-karts} for \textbf{GKDC} (Go-kart Design Challenge Competition) and \textbf{design of electric formula student vehicle} (based on the official student formula notebook) for \textbf{FGCC} (Formula Green Concept Challenge), both events organised by \textbf{ISNEE Motorsport} (Indian Society of New Era Engineers).}
		
		\resumeProjectHeading
		{\textbf{Other Projects} $|$ \emph{On-Going Projects}}{Aug'24 --}
		\begin{itemize}
			\item \small \textbf{Bi-copter:} A fully portable (can be folded and stacked together) two-rotor copter with tilt-rotor mechanism for military applications.
			\item \small \textbf{Fixed Wing VTOL:} A low-endurance tail-sitter e-UAV with single-rotor propulsion system used for surveillance. Transition and hover control with control vanes using thrust vectoring (inspired by the V-BAT).
			\item \small \textbf{Humanoid:} A full-scale anthropomorphic companion robot with VLA (Visual Language Action) integration.
			\item \small \textbf{Spider-Bot:} A 4-leg robot which can navigate in rough terrain, for space and military applications.
		\end{itemize}
		\vspace{1pt}
		
		\resumeProjectHeading
		{\textbf{Deep Koopman Based Quadrotor System Identification} $|$ \emph{Final Year Project}}{On-going}
		\vspace{-6pt}
		\resumeItemDesc{Using the simulated \textbf{flight data generated in Gazebo} to learn Koopman observables and the corresponding global linear representation in these observables space using the \textbf{deep Lipschitz neural network}. The learned system dynamics is later used to implement a \textbf{real-time linear model predictive control (LMPC)} for robust trajectory tracking under external disturbances and wind conditions. Its performance is then validated against MPC implemented on classical linear, non-linear, and other Koopman-based models of the quadrotor system.}
		
		\resumeProjectHeading
		{\textbf{Implementation of Control Strategies for a Self-Balancing Cube} $|$ \emph{Capstone Project}}{}
		\vspace{-6pt}
		\resumeItemDesc{Derived the nonlinear dynamics of a reaction-wheel self-balancing cube using \textbf{Lagrangian mechanics} and formulated the linearized state-space representation. Implemented \textbf{Linear Quadratic Regulator (LQR)} and \textbf{Active Disturbance Rejection Control (ADRC)} algorithms, demonstrating through simulation the two control strategies.}
		
		\resumeProjectHeading
		{\textbf{Sliding-Mode Control for Flight Stability of Quadrotor Drone} $|$ \emph{Adaptive Super-Twisting Reaching Law}}{}
		\vspace{-6pt}
		\resumeItemDesc{Implemented a robust \textbf{Adaptive Super-Twisting Sliding Mode Controller (ASTRL)} to enhance the flight stability of a quadrotor drone against disturbances and sensor noise. An exponential function-based adaptive law was implemented to dynamically adjust control gains, achieving \textbf{faster reaching times and reduced chattering} compared to conventional SMC and traditional Super-Twisting Algorithms.}
		
		\resumeProjectHeading
		{\textbf{Cooperative Mapping in an Unknown Grid Environment} $|$ \emph{Capstone Project}}{}
		\vspace{-6pt}
		\resumeItemDesc{Implemented \textbf{frontier-based exploration}, where robots move toward unknown cells next to explored areas. Designed a utility function that balances information gain (new cells discoverable within sensor range using line-of-sight with obstacle blocking) and travel cost (Manhattan distance). Used \textbf{Breadth-First Search (BFS)} for path planning on the partially known map and added collision avoidance so robots do not move into the same or conflicting cells. Visualized explored areas, frontiers, robot paths, and real-time utility values in a \textbf{pygame} environment.}
		
		\resumeProjectHeading
		{\textbf{Vision-Based Obstacle Detection and Avoidance for UAV Navigation} $|$ \emph{Capstone Project}}{}
		\vspace{-6pt}
		\resumeItemDesc{Demonstrated that a \textbf{minimal perception setup}, using only an \textbf{Intel RealSense} camera, is sufficient for reliable perception and real-time collision-free navigation in cluttered environments. The system includes \textbf{3D voxel-based occupancy mapping}, dynamic obstacle detection and tracking using an ensemble of \textbf{U-Depth, DBSCAN, and YOLO-MAD}, and Kalman filters for accurate state estimation. For motion planning, the \textbf{RRT} algorithm explores feasible paths, refined using \textbf{B-spline optimisation} to generate optimal, collision-free paths for quadrotors and lightweight UAVs with limited payload and power.}
		
		\resumeProjectHeading
		{\textbf{Reinforcement Learning Projects} $|$ \emph{Personal and Capstone Projects}}{}
		\begin{itemize}
			\item \small \textbf{Model-based and model-free RL:} Comparing the Adaptive Dynamic Programming (ADP) agent (model-based learning) and tabular Q-Learning agent (model-free learning) in a grid-world environment.
			\item \small \textbf{RL-based formation control:} Implemented a leader robot that follows a trajectory with a simple PID controller; the followers are independent agents that take actions by learning a policy using \textbf{Double Deep Q-Network (DDQN)}.
			\item \small \textbf{The Option-Critic Architecture:} Implemented the Option-Critic architecture to autonomously discover temporal abstractions (options). By jointly learning intra-option policies, termination functions, and a policy over options via policy gradient theorems, the agent developed high-level behaviors without manual subgoals. Results across \textbf{Fourrooms and Atari} environments showed faster adaptation to changing goals and superior transfer learning compared to primitive action methods.
		\end{itemize}
		\vspace{1pt}
		
		\resumeProjectHeading
		{\textbf{Hypothesis Testing and Exploratory Analysis of Australian Suburbs} $|$ \emph{Capstone Project}}{}
		\vspace{-6pt}
		\resumeItemDesc{Conducted \textbf{unsupervised learning} on Melbourne suburb datasets to uncover socio-economic and demographic patterns. Utilized \textbf{Principal Component Analysis (PCA)} for dimensionality reduction and engineered custom similarity metrics using Euclidean and Cosine distances. Implemented \textbf{K-Means clustering} and Multidimensional Scaling (MDS) to segment profiles, validating the hypothesis that geographical proximity correlates with profile similarity through geodesic mapping and statistical testing.}
		
		\resumeProjectHeading
		{\textbf{Earthquake-Prone Building Assessment} $|$ \emph{Building Material Classification from Street-View Imagery}}{}
		\vspace{-6pt}
		\resumeItemDesc{Developed a machine learning pipeline using \textbf{TensorFlow and Keras} to classify building structures into five material types for earthquake risk assessment from Google Street-View images. Utilized \textbf{Transfer Learning} by fine-tuning a pre-trained \textbf{ResNet50} model. Addressed class imbalance with data augmentation, and improved generalization through \textbf{L2 regularization, Dropout}, and adaptive callbacks (EarlyStopping, ReduceLROnPlateau).}
		
		\resumeProjectHeading
		{\textbf{VidQuery} $|$ \emph{Multimodal RAG System, Capstone Project}}{}
		\vspace{-6pt}
		\resumeItemDesc{A multimodal \textbf{Retrieval-Augmented Generation (RAG)} system that allows users to query educational videos through an interactive \textbf{Gradio} interface, delivering precise answers alongside exact timestamp citations. The pipeline extracts audio via \textbf{Whisper}, slide text via \textbf{PaddleOCR}, and video frames, generating text embeddings with \textbf{BGE} and visual embeddings with \textbf{SigLIP}, stored in a \textbf{Milvus} database. Agentic orchestration and multi-step retrieval use \textbf{Llama 3.1 8B Instruct}, while \textbf{Llama 3.2 11B Vision} handles deep visual frame analysis. Built entirely on an open-source stack, running locally to ensure data privacy and eliminate cloud API costs.}
		
	\end{itemize}
	
	%==================  EXPERIENCE & LEADERSHIP  ==================
	\section{Professional Experience \& Leadership}
	\begin{itemize}[leftmargin=0.15in, label={}]
		
		\resumeSubheading
		{Paraha Ventures}{Mysuru, India}
		{Co-Founder}{2022 -- 2023}
		\vspace{-6pt}
		\resumeItem{Co-founded and actively operated a \textbf{freelance digital marketing startup} for 1.5 years following graduation.}
		\resumeItem{Delivered \textbf{end-to-end digital solutions} for local businesses, managing projects across Website Development, \textbf{SEO}, and Social Media Marketing.}
		\resumeItem{Gained hands-on \textbf{entrepreneurial experience} in client relations and project management before transitioning to higher studies.}
		\vspace{-9pt}
		
		\resumeSubheading
		{Student Motorsport Team (ISNEE Motorsport)}{Mysuru, India}
		{Captain}{2020 -- 2021}
		\vspace{-9pt}
		\resumeItem{Led the team in the \textbf{design and fabrication of multiple electric karts}, managing engineering timelines and team logistics.}
		\resumeItem{Competed in prestigious national-level events including \textbf{Formula Green (Formula Student Electric)} and the \textbf{Go-kart Design Challenge (GKDC)}, both organized by ISNEE Motorsport, India, \textbf{winning several awards}.}
		\vspace{3pt}
		
		\resumeSubheading
		{Make A Difference (MAD)}{Mysuru, India}
		{Tutor}{2019 -- 2021}
		\vspace{-9pt}
		\resumeItem{Served 2 years as a tutor at \textbf{Make A Difference (MAD)}, a youth-driven NGO dedicated to supporting children in need of care and protection.}
	\end{itemize}
	
	%==================  ADDITIONAL  ==================
	\section{Additional Information}
	\begin{itemize}[leftmargin=0.15in, label={}]
		\small{\item{
				\textbf{Languages}{: Kannada, Hindi, Telugu, English} \\
				\textbf{Interests}{: Basketball, Cricket, Chess, Frisbee; Filmmaking, Writing, Character Sketching}
		}}
	\end{itemize}
	
\end{document}